% !TEX root = ../../hw01.tex
% Homework 01, Exercise 3. Shared solution.

% Transcribed from IMG_8997 and the upper left page of IMG_8998.
\begin{proof}[Solution]
We use Euclid's lemma, which states that if a prime divides a
product of integers, it divides at least one factor. Applying this
repeatedly shows that $p\mid a^k$ implies $p\mid a$ for every
positive integer $k$.

\textit{The square root.}
Suppose, for a contradiction, that $\sqrt p\in\Q$. Write
\[
  \sqrt p=\frac ab,
  \qquad a,b\in\Z,\quad b>0,\quad\gcd(a,b)=1.
\]
Squaring gives $a^2=pb^2$, so $p\mid a^2$.
Euclid's lemma gives $p\mid a$, and we can write $a=pm$
for some $m\in\Z$. Substituting and cancelling one factor of $p$ gives
\[
  p^2m^2=pb^2
  \quad\Longrightarrow\quad
  pm^2=b^2.
\]
Hence $p\mid b^2$, so Euclid's lemma also gives $p\mid b$.
The prime $p>1$ therefore divides both $a$ and $b$,
contradicting $\gcd(a,b)=1$. Thus $\sqrt p$ is irrational.

\textit{The general $k$th root.}
Let $k\ge2$ and suppose that $p^{1/k}=a/b$ in lowest terms,
with $a,b\in\Z$ and $b>0$. Raising both sides to the $k$th
power gives $a^k=pb^k$. Thus $p\mid a^k$, and Euclid's lemma
gives $p\mid a$. Writing $a=pm$, we obtain
\[
  p^km^k=pb^k
  \quad\Longrightarrow\quad
  p^{k-1}m^k=b^k.
\]
Since $k-1\ge1$, this shows that $p\mid b^k$, so $p\mid b$.
Again, $p$ divides both $a$ and $b$, contradicting that the
fraction is in lowest terms. Therefore $p^{1/k}$ is irrational
for every integer $k\ge2$.
\end{proof}
